% Part: methods
% Chapter: induction
% Section: relations

\documentclass[../../../include/open-logic-section]{subfiles}

\begin{document}

\olfileid{mth}{ind}{rel}

\olsection{العلاقات والدوال}

حين نعرّف مجموعةً من الكائنات (كالأعداد الطبيعية أو الحدود السليمة)
تعريفاً استقرائياً، يمكننا أيضاً أن نعرّف بالاستقراء \emph{علاقات على}
هذه الكائنات. انظر، مثلاً، في الفكرة الآتية: يكون الحد السليم~$t_1$
حداً فرعياً للحد السليم~$t_2$ إذا ورد جزءاً منه. ولنرمز إلى ذلك بالرمز
$t_1 \sqsubseteq t_2$. ومن البديهي أن كل حد سليم حد فرعي من نفسه:
$t \sqsubseteq t$. ويمكننا أن نعطي هذه العلاقة التعريف الاستقرائي الآتي:

\begin{defn}
  تُعرَّف بالاستقراء على~$t_2$ علاقةُ كون الحد السليم~$t_1$ حداً فرعياً
  لـ~$t_2$، أي $t_1 \sqsubseteq t_2$، كما يأتي:
  \begin{enumerate}
  \item إذا كان $t_2$ حرفاً، فإن $t_1 \sqsubseteq t_2$ إذا وفقط إذا كان
    $t_1 = t_2$.
  \item إذا كان $t_2$ هو $[s_1 \circ s_2]$، فإن
    $t_1 \sqsubseteq t_2$ إذا وفقط إذا كان $t_1 = t_2$، أو
    $t_1 \sqsubseteq s_1$، أو $t_1 \sqsubseteq s_2$.
  \end{enumerate}
\end{defn}

يخبرنا هذا التعريف، مثلاً، بأن $
\mathrm{a} \sqsubseteq [\mathrm{b} \circ \mathrm{a}]$. ذلك أن البند~(2)
يقول إن $
\mathrm{a} \sqsubseteq [\mathrm{b} \circ \mathrm{a}]$ إذا وفقط إذا كان
$\mathrm{a} = [\mathrm{b} \circ \mathrm{a}]$، أو
$\mathrm{a} \sqsubseteq \mathrm{b}$، أو
$\mathrm{a} \sqsubseteq \mathrm{a}$. والشرطان الأولان كاذبان: فمن الواضح
أن $\mathrm{a}$ لا يطابق $[\mathrm{b} \circ \mathrm{a}]$، وبحسب~(1)،
فإن $\mathrm{a} \sqsubseteq \mathrm{b}$ إذا وفقط إذا كان
$\mathrm{a} = \mathrm{b}$، وهذا كاذب أيضاً. أما، وبحسب~(1) كذلك،
$\mathrm{a} \sqsubseteq \mathrm{a}$ إذا وفقط إذا كان
$\mathrm{a} = \mathrm{a}$، وهذا صادق.

من المهم ملاحظة أن نجاح هذا التعريف يتوقف على حقيقة لم نبرهنها بعد:
كل حد سليم~$t$ إما أن يكون حرفاً منفرداً، وإما أن توجد حدود سليمة
$s_1$ و$s_2$ \emph{محددة على نحو فريد} بحيث
$t = [s_1 \circ s_2]$. والمقصود بعبارة «محددة على نحو فريد» هنا أنه إذا
كان $t = [s_1 \circ s_2]$، فلا يكون \emph{أيضاً}
$= [r_1 \circ r_2]$ مع $s_1 \neq r_1$ أو $s_2 \neq r_2$.
ولو وقع ذلك، لأمكن أن يتعارض البند~(2) مع نفسه: فقد نحصل، عند قراءة
$t_2$ على أنه $[s_1 \circ s_2]$، على $t_1 \sqsubseteq t_2$، لكننا قد
نحصل، عند قراءته على أنه $[r_1 \circ r_2]$، على نفي
$t_1 \sqsubseteq t_2$. وقبل أن نبرهن استحالة وقوع ذلك، فلننظر في مثال
يمكن أن يقع فيه بالفعل.

\begin{defn}
  نعرّف \emph{الحدود الخالية من الأقواس} استقرائياً كما يأتي:
  \begin{enumerate}
  \item كل حرف حد خالٍ من الأقواس.
  \item إذا كان $s_1$ و$s_2$ حدين خاليين من الأقواس، فإن
    $s_1 \circ s_2$ حد خالٍ من الأقواس.
  \item لا شيء آخر حد خالٍ من الأقواس.
  \end{enumerate}
\end{defn}

من أمثلة الحدود الخالية من الأقواس: $\mathrm{a}$، و
$\mathrm{b} \circ \mathrm{d}$، و
$\mathrm{b} \circ \mathrm{a} \circ \mathrm{b}$. فإذا عرّفنا «الحد
الفرعي» للحدود الخالية من الأقواس بالطريقة السابقة نفسها، كان نص البند
الثاني:
\begin{center}
  إذا كان $t_2 = s_1 \circ s_2$، فإن $t_1 \sqsubseteq t_2$ إذا وفقط إذا
  كان $t_1 = t_2$، أو $t_1 \sqsubseteq s_1$، أو
  $t_1 \sqsubseteq s_2$.
\end{center}

والآن فإن $\mathrm{b} \circ \mathrm{a} \circ \mathrm{b}$ من الصورة
$s_1 \circ s_2$، حيث
\begin{align*}
  s_1 & = \mathrm{b} \text{ و} & s_2 & = \mathrm{a} \circ \mathrm{b}.
\intertext{وهو أيضاً من الصورة $r_1 \circ r_2$، حيث}
  r_1 & = \mathrm{b} \circ \mathrm{a} \text{ و} & r_2 &= \mathrm{b}.
\end{align*}
فهل $\mathrm{a} \circ \mathrm{b}$ حد فرعي من
$\mathrm{b} \circ \mathrm{a} \circ \mathrm{b}$؟ يكون الجواب نعم إذا
أخذنا بالقراءة الأولى، ولا إذا أخذنا بالقراءة الثانية.

تسمى خاصيةُ كون الطريقة التي يُبنى بها الحد السليم من حدود سليمة أخرى
فريدةً \emph{قابليةَ القراءة الفريدة}. ولأن التعريفات الاستقرائية
للعلاقات على مثل هذه الكائنات المعرفة استقرائياً مهمة، فعلينا أن نبرهن
تحقق هذه الخاصية.

\begin{prop}
  لنفترض أن $t$ حد سليم. عندئذ إما أن يكون $t$ حرفاً منفرداً، وإما أن
  يوجد حدان سليمان محددان على نحو فريد، $s_1$ و$s_2$، بحيث
  $t = [s_1 \circ s_2]$.
\end{prop}

\begin{proof}
  إذا كان $t$ حرفاً منفرداً، فالشرط متحقق. فلنفترض إذن أن $t$ ليس حرفاً
  منفرداً. نستدل من التعريف الاستقرائي على أنه لا بد أن يكون من الصورة
  $[s_1 \circ s_2]$ لبعض الحدين السليمين $s_1$ و~$s_2$. ويبقى أن نبين
  أنهما محددان على نحو فريد؛ أي إنه إذا كان $t = [r_1 \circ r_2]$، فإن
  $s_1 = r_1$ و$s_2 = r_2$.

  فلنفترض أن $t = [s_1 \circ s_2]$ وأن $t = [r_1 \circ r_2]$ أيضاً،
  حيث $s_1$ و$s_2$ و$r_1$ و$r_2$ حدود سليمة. علينا أن نبين أن
  $s_1 = r_1$ و$s_2 = r_2$. أولاً، لا بد أن يتطابق $s_1$ و$r_1$، وإلا
  لكان أحدهما مقطعاً ابتدائياً حقيقياً للآخر. لكن ذلك، بموجب
  \olref[sti]{prop:initial}، مستحيل إذا كان كل من $s_1$ و~$r_1$ حداً
  سليماً. وإذا كان $s_1 = r_1$، فمن الواضح أن $s_2 = r_2$ أيضاً.
\end{proof}

ويمكننا كذلك تعريف الدوال استقرائياً؛ فمثلاً، يمكننا تعريف الدالة~$f$
التي تسند إلى كل حد سليم العمق الأقصى للأقواس المتداخلة $[\dots]$ فيه كما
يأتي:

\begin{defn}
  \ollabel{defn:depth} يُعرَّف \emph{عمق} الحد السليم، $f(t)$، استقرائياً
  كما يأتي:
  \[
  f(t) = \begin{cases}
    0 & \text{ إذا كان $t$ حرفاً}\\
    \max(f(s_1), f(s_2)) + 1 & \text{ إذا كان $t = [s_1 \circ s_2]$.}
  \end{cases}
  \]
\end{defn}

فعلى سبيل المثال:
\begin{align*}
  f([\mathrm{a} \circ \mathrm{b}]) & =
    \max(f(\mathrm{a}),f(\mathrm{b})) + 1 \\
  &= \max(0, 0) + 1 = 1 \text{، و}\\
  f([[\mathrm{a} \circ \mathrm{b}] \circ \mathrm{c}]) & =
    \max(f([\mathrm{a} \circ \mathrm{b}]), f(\mathrm{c})) + 1 \\
  & = \max(1,0) + 1 = 2.
\end{align*}

نفترض هنا، بطبيعة الحال، أن $s_1$ و$s_2$ حدان سليمان، ونستعمل حقيقة أن
كل حد سليم إما حرف وإما من الصورة $[s_1 \circ s_2]$. ومن المهم، مرة
أخرى، أنه لا يمكن أن يكون من هذه الصورة إلا بطريقة واحدة. ولرؤية السبب،
لنعد إلى الحدود الخالية من الأقواس التي عرّفناها آنفاً. سيكون «التعريف»
المقابل:
\[
  g(t) =
  \begin{cases}
    0 & \text{ إذا كان $t$ حرفاً}\\
   \max(g(s_1), g(s_2)) + 1 & \text{ إذا كان $t = s_1 \circ s_2$.}
  \end{cases}
\]
ولننظر الآن في الحد الخالي من الأقواس
$\mathrm{a} \circ \mathrm{b} \circ \mathrm{c} \circ \mathrm{d}$.
يمكن قراءته بأكثر من طريقة؛ منها، مثلاً، أن يُقرأ على أنه
$s_1 \circ s_2$، حيث
\begin{align*}
  s_1 & = \mathrm{a} \text{ و} &
  s_2 & = \mathrm{b} \circ \mathrm{c} \circ \mathrm{d},
\intertext{أو على أنه $r_1 \circ r_2$، حيث}
  r_1 & = \mathrm{a} \circ \mathrm{b} \text{ و} &
  r_2 &= \mathrm{c} \circ \mathrm{d}.
\end{align*}
يعطي حساب $g$ بحسب القراءة الأولى:
\begin{align*}
  g(s_1 \circ s_2) & =
  \max(g(\mathrm{a}), g(\mathrm{b} \circ \mathrm{c} \circ \mathrm{d})) + 1 \\ & =
  \max(0,2) + 1 = 3
  \intertext{بينما نحصل بحسب القراءة الأخرى على}
  g(r_1 \circ r_2) & =
  \max(g(\mathrm{a} \circ \mathrm{b}), g(\mathrm{c} \circ \mathrm{d})) + 1 \\ &
  = \max(1,1) + 1 = 2.
\end{align*}
لكن لا بد أن تعطي الدالة دائماً قيمة فريدة؛ ولذلك فإن «تعريفنا» لـ~$g$
لا يعرّف دالة أصلاً.

\begin{prob}
  أعط تعريفاً استقرائياً للدالة $l$، حيث $l(t)$ هو عدد الرموز في الحد
  السليم~$t$.
\end{prob}

\begin{prob}
  برهن بالاستقراء البنيوي على الحدود السليمة $t$ أن $f(t) < l(t)$ (حيث
  $l(t)$ عدد الرموز في~$t$، و$f(t)$ عمق $t$ كما عُرّف في
  \olref[mth][ind][rel]{defn:depth}).
\end{prob}

\end{document}
